\worksheettitle
  {M07-R. Считаем переходы}
  {\modulemark{M07-R} Коррекционная карточка}

\studentnameblock

\begin{warningbox}[Назовите ошибку]
  Расстояние составлено из промежутков между объектами, а не из самих
  объектов. Сначала покажите переходы на коротком ряду.
\end{warningbox}

\equalpointsbasic

\begin{practicebox}[Шаг 1. Посчитайте по рисунку]
  На рисунке \answerblank[15mm] точек и \answerblank[15mm] промежутков.
  Между 2-й и 6-й точками \answerblank[15mm] промежутка.
\end{practicebox}

\begin{practicebox}[Шаг 2. Перейдите к номерам]
  \begin{enumerate}[label=\alph*)]
    \item между 4-й и 9-й: $9-4=\answerblank[15mm]$;
    \item между 12-й и 20-й: \answerblank[15mm];
    \item между соседними 35-й и 36-й: \answerblank[15mm].
  \end{enumerate}
\end{practicebox}

\begin{practicebox}[Шаг 3. Добавьте длину]
  Между 3-й и 8-й точками шаг равен $4$ см.
  \[
    \text{промежутков } \answerblank[15mm],\qquad
    \text{расстояние } \answerblank[15mm]\cdot4
      =\answerblank[18mm]\text{ см}.
  \]
\end{practicebox}

\begin{practicebox}[Шаг 4. Исправьте]
  Ученик написал: «Шесть столбов через $3$ м занимают $6\cdot3=18$ м».

  Причина: \answerblank[95mm]

  Исправление: \answerblank[95mm]
\end{practicebox}

\begin{checkpointbox}[Повторная проверка]
  Восемь деревьев стоят через $5$ м, включая оба крайних.
  Найдите число промежутков и расстояние между крайними деревьями.
  Объясните оба действия. \writinglines{3}
\end{checkpointbox}
